\section{Data preparation}
\label{sec:prep}

\Cref{tab:prep} lists every preparation decision. All of them are implemented in
\tcode{src/data\_prep.py} and in the scikit-learn pipeline.

\begin{table}[!htbp]
  \centering
  \caption{Preparation decisions, their justification, and what each one risks.}
  \label{tab:prep}
  \footnotesize
  \begin{tabularx}{\linewidth}{>{\raggedright\arraybackslash}p{3.6cm} X}
    \toprule
    \textbf{Decision} & \textbf{Why, and what it risks} \\
    \midrule
    Drop 7{,}754 cancelled or suspected-fraud lines &
      These orders are never delivered, so there is no delivery decision for them. We assume
      the cancellation or fraud check happens before dispatch. \\
    Aggregate lines to orders &
      The decision is per order. Keeping lines would count multi-line orders several times and
      let lines of one order fall on both sides of a CV split. \\
    Order-level features &
      Lines, quantity, value, discount, maximum unit price, number of categories, profit,
      margin, and the department and category of the highest-value line. \\
    Calendar features &
      Weekday, hour and month of the order. The raw timestamp and \tcode{Order Id} are
      \emph{not} used, because they encode the order sequence, and the sequence determines the
      outcome here (\Cref{tab:realism}). \\
    History features &
      Prior order count and smoothed prior late rate, $(\text{late}+2.5)/(n+5)$, per customer,
      city, and city $\times$ mode. A past order counts only once its outcome was known. \\
    Drop PII and empty columns &
      Email, password, names, street and zip codes; \tcode{Product Description},
      \tcode{Order Zipcode} and \tcode{Product Image}. \\
    Keep outliers &
      They are genuine prices and losses, and trees are insensitive to them. \\
    Encoding and scaling &
      One-hot encoding for 6 low-cardinality categoricals (levels with fewer than 20 orders are
      pooled). Target encoding\cite{micci2001}, cross-fitted over 5 folds, for country (164),
      city (3{,}585) and category (50). Median imputation. Standard scaling for non-tree
      models only. \\
    \bottomrule
  \end{tabularx}
\end{table}

\paragraph{Leakage risks specific to this dataset and how each is avoided.}
\begin{enumerate}
  \item \textbf{\tcode{shipping date (DateOrders)}} looks like a dispatch date. In fact it equals the
        order date plus the \emph{real} transit days in 97.4\% of rows, so it encodes the outcome.
        It is excluded, together with \tcode{Days for shipping (real)}, \tcode{Delivery Status}
        and \tcode{Order Status}.
  \item \textbf{Order sequence.} Transit days follow \tcode{Order Id} modulo 5, so a model given
        the order ID or the raw timestamp could learn how the data was generated. Neither is a
        feature, and \tcode{run\_all.py} stops if one is ever added.
  \item \textbf{History features.} A naive ``past late rate'' would include orders that had not
        yet been delivered. We count a past order only after its order date + real transit
        days + 1 day.
  \item \textbf{Preprocessing.} Imputers, scalers and encoders sit inside the \tcode{Pipeline} and
        are refit on the training part of every fold. Target encoding is also cross-fitted
        within the training data.
\end{enumerate}
